\documentclass[10pt,a4paper]{amsart}
\usepackage[margin=19mm]{geometry}
\usepackage{amsmath,amssymb,mathtools}
\usepackage[scr=rsfs,cal=euler]{mathalfa}
\usepackage{xfrac}
\usepackage[unicode,hidelinks,bookmarks=false]{hyperref}
\newcommand{\ie}{i.e.\ }
\newcommand{\cf}{cf.\ }
\newcommand{\resp}{respectively\ }
\newcommand{\Subsection}{\subsection}
\providecommand{\nobreakdash}{\nobreak\hbox{-}}
\setlength{\emergencystretch}{2em}
\multlinegap0pt
\usepackage{amssymb}
\usepackage{mathtools}
\newtheorem{thm}{Theorem}
\newtheorem{lem}{Lemma}
\def\R{\mathbb R}
\title{A simple integral representation of single-event scoring rules}
\author{Alexander R. Pruss}
\date{Source: arXiv:2603.00750v1 (CC BY 4.0)}
\begin{document}
\maketitle

\noindent\emph{Source and licence.} A simple integral representation of single-event scoring rules. Version arXiv:2603.00750v1 (CC BY 4.0), distributed by arXiv under the Creative Commons Attribution 4.0 International licence, http://creativecommons.org/licenses/by/4.0/, as recorded in the arXiv licence metadata for this exact version and shown on the abstract page https://arxiv.org/abs/2603.00750v1. Full licence notice: \texttt{licenses/arxiv-2603.00750-CC-BY-4.0.txt}.\par\smallskip

\noindent\textit{Editorial scope.} Counted unit: the article's thm thm:repr (source0.tex lines 167--178). The article's complete original proof of it is delivered with it (source0.tex lines 266--372, one \texttt{proof} environment of 107 lines, which the article places away from the statement). No mathematics is written, repaired or re-derived: the statement and the proof are the article's own lines, in the article's own notation, and no retained line is substituted or re-flowed.  Retained uncounted as the source-local context the proof needs: lines 210--212.  Nothing was omitted from any retained span, so the package carries no omission record.  No retained line contains a citation, so the wrapper carries no bibliography and no bibliography entry.  No retained line contains a figure environment, an image-inclusion command or drawing code, so no image is delivered and the package declares no asset dependency.  Editorial formatting only, in addition: an AMS \texttt{amsart} preamble that restates the article's own declarations and macros used by the retained lines, namely the environments \texttt{thm}, \texttt{lem} on separate counters, together with the standard packages those lines need.  The retained environments print in this document as Theorem 1, Lemma 1 in order of appearance, whereas the article prints the same items with the numbering of its own full document.  The complete native source of the article as distributed in the arXiv e-print for this exact version is bundled unchanged as \texttt{sources/195506/source0.tex}.
\begin{lem}\label{lem:unique}
Suppose $(T,F)$ and $(T,K)$ are both proper. Then $F$ and $K$ differ by a constant on $[0,1)$.
\end{lem}

\begin{thm}\label{thm:repr}
Let $T:[0,1]\to[-\infty,\infty)$ be monotonic non-decreasing and finite on $(0,1)$.
Fix any real constant $C$ and any $c\ge 0$. Let
$$
	F(x) = C-\frac{xT(x)}{1-x}+\int_{1/2}^x \frac{T(u)}{(1-u)^2} \, du,
$$
for $x<1$. If $T$ is continuous at $1$, let $F(1)=-c+\lim_{x\to 1-} F(x)$, and otherwise
let $F(1)=-\infty$. Then the limit is defined if $T$ is continuous at $1$, and in any case $(T,F)$ is proper.

Conversely, any function $F$ on $[0,1]$ that is finite except perhaps at $1$ and such that
$(T,F)$ is proper satisfies the above definition for some real $C$ and $c\ge 0$. 
\end{thm}

\begin{proof}[Proof of Theorem~\ref{thm:repr}]
First consider the case where $T(x)$ is an indicator function of an interval $A \subseteq (0,1]$ containing
$1$, where $1_A(x)$ is $1$ if
$x\in A$ and $0$ if $x\notin A$. Let:
$$
	G_A(x) = -\frac{xT(x)}{1-x}+\int_0^x \frac{T(u)}{(1-u)^2} \, du
$$
for $x<1$. For $x \in (0,1)-A$ we then have $F_A(x)=0$. For $x \in A\cap (0,1)$, we have $x\ge a$ and:
$$
\begin{aligned}
	G_A(x) &= -\frac{x}{1-x}+\int_a^x (1-u)^{-2} \,du \\
		 &= -\frac{x}{1-x}+\frac{1}{1-x}-\frac{1}{1-a} \\
		 &= -\frac{a}{1-a}.
\end{aligned}
$$		 
Thus $G_A(x) = -(a/(1-a))\cdot 1_A(x)$ for $x<1$. If $A=\{ 1 \}$, then $F_A$ is discontinuous at $1$,
and let $G_A(1)=-\infty$. Otherwise, $F_A$ is continuous at $1$, and let $-G_A(1)=(a/(1-a))\cdot 1_A(1)$, 
with the convention that this is everywhere $0$ if $A=\varnothing$. It is easy to check that 
$(1_A,G_A)$ is proper.

Now let $T_B = -1_B$ where $B$ is of the form $[0,b]$ and $[0,b)$ for $b\le 1$. Let:
$$
	F_B(x) = -\frac{xT(x)}{1-x}+\int_{1/2}^x \frac{T(u)}{(1-u)^2} \, du
$$
for $x<1$. Let $F_B(1)=-\infty$ if $B=[0,1)$ and otherwise let $F_B(1)=\lim_{x\to 1-} F_B(x)$.
If $b<1$, then $T_B$ differs by a constant from $1_{B^c}$ and $F_B$ by a constant from $G_{B^c}$. Thus
$(T_B,F_B)$ is also proper by what we proved already. If $B=[0,1]$, then $T_B$ and $F_B$ are both constant, and
$(T_B,F_B)$ are proper. If $B=[0,1)$, then $F_B$ is $-\infty$ at $1$ and on $[0,1)$ it's constant. In that
case it is easy to verify the propriety inequality
$$
	p T_B(p) + (1-p) F_B(p) \ge p T_B(x) + (1-p) F_B(x).
$$
For both sides are equal if $p,x<1$. If $p=1$, the inequality is equivalent to $p T_B(p) \ge p T_B(x)$ which
holds by monotonicity. If $x=1$ and $p<1$, the right-hand-side is $-\infty$. So, in all cases we have
$(T_B,F_B)$ proper.

Now given any non-decreasing $T(x)$ bounded below, without loss of generality assume $T$ is non-positive
(subtract a constant if necessary). Let $B_t=\{ x : T(x) \le t \}$ be a level set of $T$. Then $B_t$ is of the form 
$[0,b]$ or $[0,b)$ and:
\begin{equation}\label{eq:Tint}
	T(x) = \int_{-\infty}^0 T_{B_t} \, dt.
\end{equation}
Let $F$ be as in the statement of the Lemma in the special case where $C=c=0$. Thus:
$$
	F(x) = -\frac{xT(x)}{1-x}+\int_{1/2}^x \frac{T(u)}{(1-u)^2} \, du,
$$
and $F(1)$ is $-\infty$ if $T$ is continuous at $1$ and is $\lim_{x\to1-} F(x)$ otherwise.
We will show that $F(T,F)$ is proper.

For $x<1$ we have:
\begin{equation}\label{eq:Fint}
\begin{aligned}
	F(x) &= -\frac{xT(x)}{1-x}+\int_{1/2}^x \frac{T(u)}{(1-u)^2} \, du \\
		 &= -\frac{x}{1-x}\int_{-\infty}^0 T_{B_t} \, dt + \int_{1/2}^x \int_{-\infty}^0 \frac{T_{B_t}}{(1-u)^2} \, dt\, du \\
		 &= \int_{-\infty}^0 \left(-\frac{xT_{B_t}}{1-x}  + \int_{1/2}^x \frac{T_{B_t}}{(1-u)^2} \, du \right) \, dt\\
		 &= \int_{-\infty}^0 F_{B_t}(x) \, dt,
\end{aligned}		 
\end{equation}
where the order of integration was swapped by Tonelli's theorem as $T\le 0$. By propriety of
$(T_{B_t},F_{B_t})$ we have:
$$
	p T_{B_t}(p) + (1-p) F_{B_t}(p) \ge 
	p T_{B_t}(x) + (1-p) F_{B_t}(x),
$$
for $p,x<1$, 
and integrating from $-\infty$ to $0$ with respect to $t$, by \eqref{eq:Tint} and \eqref{eq:Fint} we get
\begin{equation}\label{eq:propriety}
	p T(p) + (1-p) F(p) \ge 
	p T(x) + (1-p) F(x).
\end{equation}

Note that $F_{B_t}$ differs by a constant from $G_{B_t^c}$ on $[0,1)$, and hence is non-increasing on $[0,1)$.
By \eqref{eq:Fint} this is also true of $F$, so it has a limit at $1$. Define $F$ to be that limit if
$T$ is continuous at $1$ and to be $-\infty$ otherwise. In the case where $T$ is continuous at $1$, we get
\eqref{eq:propriety} for $p=1$ and $x<1$ by taking the limit as $p\to 1-$, and for $p<1$ and $x=1$ by
taking the limit as $x\to 1-$. The case of $p=x=1$ is trivial. What remains is the case where $T$ is
not continuous at $1$ and either $p=1$ or $x=1$. If $p=1$, then \eqref{eq:propriety} follows
from the monotonicity of $T$, and if $p<1$ but $x=1$, then \eqref{eq:propriety} is trivial when
$F(1)=-\infty$.

Now fix $C\in\R$ and $c\ge 0$. Let $K(x)=F(x)+C$ for $x<1$ and $K(1)=F(1)+C-c$. I claim that $(F,K)$ is proper.
For \eqref{eq:propriety}
remains true with $K$ in place of $F$ if $p$ and $x$ are in $[0,1)$. If $p=1$, then $(1-p)K(p)=(1-p)F(p)=0$,
and so \eqref{eq:propriety} is also true with $K$ in place of $F$. The remaining case to check where $p<1$ and
$x=1$. By \eqref{eq:propriety} we have:
$$
	p T(p) + (1-p) F(p) \ge p T(1) + (1-p) F(1).
$$
But $F(p)=K(p)-C$ and $F(1)\ge K(1)-C$, so we can indeed put $K$ in place of $F$.

It remains to show that for any $K$ such that $(T,K)$ is proper and $K$ is finite except perhaps at $1$
there are $C$ and $c$ such that $K(x)=F(x)+C$ for $x<1$ and $K(1)=F(1)+C-c$. The $x<1$ part is 
Lemma~\ref{lem:unique}. It remains to show that $K(1)\le F(1)+C$. Suppose first that $T$ is discontinuous
at $1$. By propriety of $(T,K)$, for $p<1$ we have:
$$
	p (T(p)-T(1) \ge (1-p)(K(1)-K(p)).
$$
Thus,
$$
	K(1)-K(p) \le \frac{p}{1-p} (T(p)-T(1)).
$$
Since $T$ is non-decreasing and discontinuous at $1$, the limit of the right hand side as $p\to 1-$ is
$-\infty$, and hence $K(1)=-\infty=F(1)+C$. 

Now suppose $T$ is continuous at $1$, so $K(1)=F(1)+C=C+\lim_{p\to 1-} F(p)$. But
$K(1)\le \lim_{p\to 1-}$ by monotonicity, so $K(1)=F(1)+C-c$ for some positive $c$.
\end{proof}

\end{document}
